\begin{claim}{F.8}{Proved}{The seven-point kernel and its peripheral restrictions}
In the regular $A=A_4$ labelling of (F.7.8), all twelve
seven-point subsets in (F.7.9) form one free orbit for
left multiplication by $A$.  Consequently their kernels
on the specified $F_{13}$ are equal; their kernels on
its profinite completion are equal as well.
The completed kernel meets every conjugate of every
full source cusp subgroup $\widehat\Z$ trivially.
The discrete kernel contains no nonidentity reducible
mapping class under the specified index-twelve
inclusion $F_{13}\le P\Gamma_{0,4}$.
\end{claim}
\begin{proof}

Here is the actual comparison of maps.  Regard the target
as $M_{0,I}$ with its seven labels $I\subset A$, and write
\[
 f_I:X\longrightarrow M_{0,I},\qquad
             z\longmapsto[(gz)_{g\in I}].
\]
For $h\in A$, let
$r_h:M_{0,I}\xrightarrow{\sim}M_{0,hI}$ be the relabelling
which assigns to the label $hg$ the former coordinate
with label $g$.  The configuration defining $f_{hI}(z)$
has coordinate $h(gz)$ at that label.  Applying the same
M\"obius transformation $h^{-1}$ to every coordinate
therefore gives the exact identity of maps
\[
                         f_{hI}=r_h\circ f_I.
\]
This identity also identifies the base configurations;
it induces an isomorphism of the target fundamental
groups intertwining the two homomorphisms from
$\pi_1(X)$.  Functoriality of profinite completion gives
the same commuting diagram for their completions.
Since $r_h$ and its completion are isomorphisms, the
two kernel equalities follow.

The exact label certificate is
\path{code/part_a/seven_point_same_kernel.py}.
It constructs the twelve subsets in the action
$A_5/C_5$ and gives all twelve left multipliers and
the corresponding seven-label permutations.  With
the certificate's zero-based vertex numbering, one
subset is $I=\{0,1,2,3,4,7,10\}$.  Its intersections
with the three cosets of the normal subgroup
$V_4\le A_4$ are respectively
\[
 \begin{array}{c|c}
 \text{coset}&I\text{ retained in that coset}\\ \hline
 \{0,2,7,11\}&\{0,2,7\}\\
 \{1,4,5,9\}&\{1,4\}\\
 \{3,6,8,10\}&\{3,10\}.
 \end{array}
\]
The program checks that none of the twelve subsets
contains a full $V_4$ coset.  Thus the four-label
$V_4$-orbit cross-ratio is not one of the forgetful
projections of any of these seven-point maps.

The intersection of the twelve kernels in (F.7.9)
is therefore the kernel of any one of these named
maps.  F.9 below proves it nontrivial by an exact braid certificate.
This discrete kernel contains no nonidentity reducible
mapping class in the source subgroup of
$P\Gamma_{0,4}$.  To prove this, use the exact
$A_4$ covering
\[
 X=S^2\setminus R_A\longrightarrow
 S^2\setminus\{b_2,b_3,b'_3\}
\]
and its index-twelve injection on fundamental groups
in (F.7.8).  On the four-punctured quotient sphere,
every essential curve separates the moving puncture
$q$ and one branch point $b$ from the other two
branch points.  Let $e=2$ or $3$ be the order at $b$.
The preimage of the disk containing $q,b$ consists
of $12/e$ disks, each containing $e$ of the twelve
points above $q$.  These assertions hold for every
such curve: its monodromy subgroup is the conjugate
of the inertia subgroup at $b$, so its connected
components have degree $e$, and each is a disk
branched only at the one point above $b$.

The lift of a nonzero twist power belonging to the
index-twelve subgroup is a product of twists about
these boundary curves.  More precisely, its exponent
is divisible by $e$, and each lifted twist has exponent
equal to the original exponent divided by $e$.  After
forgetting five points, at least one of these boundaries
remains essential: seven retained points distributed
among six pairs or four triples force a disk containing
at least two retained points, and its complement
contains at least four.  The surviving twists all have
the same nonzero sign; parallel essential components
add their exponents.  Independence of twists about
pairwise disjoint nonisotopic essential curves makes
the resulting mapping class nonidentity.  Finally,
each nonidentity reducible pure mapping class of a
four-punctured sphere is a nonzero power of one such
twist: its invariant multicurve has just one component,
and cutting along it gives two three-punctured spheres.
This proves the asserted exclusion of reducible
elements from the discrete kernel.

There is also an injective detector for each full
profinite cusp.  Let $c$ be a primitive peripheral
generator at any one of the fourteen punctures of
$X$.  Its image in $P\Gamma_{0,4}$ is a conjugate
of $T^e$, where $T$ is a twist as above.  Thus its
lift is the indicated product of twists of exponent
one.  Choose two retained labels in an essential
lifted disk and two outside that disk.  Forgetting
all other labels gives a specified map to
$P\Gamma_{0,4}$.  The inside disk contributes one
twist $U$.  Among the other disks, another twist
survives exactly when both outside labels lie in
one disk; it then gives the same isotopy class and
the same sign.  Hence the image of $c$ is $U^d$
with $d\in\{1,2\}$, up to conjugation.

There is a character
$\ell:P\Gamma_{0,4}\to\Z$ with $\ell(U)=1$.
For example the three peripheral twists in the
free group $P\Gamma_{0,4}=F_2$ have abelianization
vectors $(1,0),(0,1),(-1,-1)$, and an essential
twist is conjugate to one of these under the pure
group.  Choose the corresponding integral linear
functional.  Completing this character and the
two forgetting maps gives a continuous homomorphism
\[
 \widehat F_{13}\longrightarrow
 \widehat{P\Gamma}_{0,7}\longrightarrow
 \widehat{P\Gamma}_{0,4}\longrightarrow\widehat\Z
\]
whose restriction to $\overline{\langle c\rangle}$
is $a\mapsto da$.  Multiplication by $1$ or $2$
is injective on $\widehat\Z$.  The cusp is a full
embedded $\widehat\Z$ in $\widehat F_{13}$, since
its primitive generator is a member of a free
basis after choosing the omitted sphere peripheral.
The same final character detects every conjugate
$g c^a g^{-1}$ with $g\in\widehat F_{13}$, because
its target is abelian.  This proves the asserted
trivial intersection of the completed kernel with
all conjugate cusp groups.  The program gives the
four retained labels and $d$ in each of the
$12\cdot7=84$ subset--inertia-partition checks;
the two fixed points of each inertia subgroup
have the same partition.

The exclusion just proved confines nonidentity discrete
kernel elements to pseudo-Anosov mapping classes. F.9 supplies
an explicit such element by an exact braid certificate.
No assertion about a Nielsen--Thurston type for profinite
elements is used.

The preceding rank-six calculation is a calculation
of its induced abelian map, and does not identify
the kernel of the fundamental-group homomorphism.
The separate file
\path{code/part_a/seven_point_braid_explore.py}
contains a numerical crossing exploration; it is
explicitly exploratory and is not used in this
map comparison. F.9 replaces the discovery computation
by a separate exact rational certificate of a kernel relation.
\end{proof}
